\documentclass{ctexart}
\usepackage{amsmath, amssymb}
\usepackage{geometry}
\geometry{a4paper, margin=2cm}

\title{\textbf{第10章 恒定磁场} --- 知识点与公式}
\author{}
\date{}

\begin{document}
\maketitle

\section*{10.1 稳恒电流}

电流强度 $I = \dfrac{\mathrm{d}Q}{\mathrm{d}t}$

电流密度 $\vec{j}$，$I = \displaystyle\int_S \vec{j}\cdot\mathrm{d}\vec{S}$

欧姆定律微分形式：$\vec{j} = \sigma\vec{E}$

电流连续性方程：$\displaystyle\oint_S \vec{j}\cdot\mathrm{d}\vec{S} = -\dfrac{\mathrm{d}}{\mathrm{d}t}\int_V \rho\,\mathrm{d}V$，稳恒时 $\nabla\cdot\vec{j} = 0$

\section*{10.2 毕奥--萨伐尔定律}

\[
\mathrm{d}\vec{B} = \frac{\mu_0}{4\pi}\frac{I\mathrm{d}\vec{l}\times\hat{r}^{0}}{r^{2}},\quad \mu_0 = 4\pi\times10^{-7}\,\text{N/A}^{2}
\]

载流直导线：$B = \dfrac{\mu_0 I}{4\pi a}(\cos\theta_1 - \cos\theta_2)$，无限长：$B = \dfrac{\mu_0 I}{2\pi a}$

载流圆线圈轴线：$B = \dfrac{\mu_0 I R^{2}}{2(R^{2}+x^{2})^{3/2}}$，圆心：$B = \dfrac{\mu_0 I}{2R}$

圆弧圆心：$B = \dfrac{\mu_0 I\phi}{4\pi R}$

螺线管轴线：$B = \dfrac{\mu_0 nI}{2}(\cos\beta_2 - \cos\beta_1)$，无限长：$B = \mu_0 nI$

运动电荷磁场：$\vec{B} = \dfrac{\mu_0}{4\pi}\dfrac{q\vec{v}\times\hat{r}^{0}}{r^{2}}$

\section*{10.3 磁高斯定理}

磁通量：$\Phi_m = \displaystyle\int_S \vec{B}\cdot\mathrm{d}\vec{S}$

磁高斯定理：$\displaystyle\oint_S \vec{B}\cdot\mathrm{d}\vec{S} = 0$（磁场无源）

\section*{10.4 安培环路定理}

$\displaystyle\oint_L \vec{B}\cdot\mathrm{d}\vec{l} = \mu_0\sum I_i$（内）

无限长圆柱均匀电流：$r > R: B = \dfrac{\mu_0 I}{2\pi r}$，$r < R: B = \dfrac{\mu_0 Ir}{2\pi R^{2}}$

无限大平面电流：$B = \dfrac{\mu_0 i}{2}$（$i$ 为面电流密度）

\section*{10.5 磁场对电流的作用}

安培力：$\mathrm{d}\vec{F} = I\mathrm{d}\vec{l}\times\vec{B}$

均匀磁场中弯曲导线：$\vec{F} = I\vec{L}_{\perp}\times\vec{B}$（$L_{\perp}$ 为起终点连线）

磁力矩：$\vec{M} = \vec{p}_m\times\vec{B}$，磁矩 $\vec{p}_m = IS\vec{n}$

磁力矩做功：$A = I\Delta\Phi_m$

\section*{10.6 带电粒子在磁场中的运动}

洛伦兹力：$\vec{f} = q\vec{v}\times\vec{B}$

回旋半径：$R = \dfrac{mv}{qB}$，周期：$T = \dfrac{2\pi m}{qB}$

螺旋运动螺距：$h = v_{\parallel}T = \dfrac{2\pi mv\cos\theta}{qB}$

霍尔效应：$U_H = \dfrac{IB}{nqd} = K\dfrac{IB}{d}$

\section*{10.7 物质的磁性}

磁介质分类：顺磁质 $\mu_r > 1$，抗磁质 $\mu_r < 1$，铁磁质 $\mu_r \gg 1$

磁场强度：$\vec{H} = \dfrac{\vec{B}}{\mu_0} - \vec{M}$，$\vec{B} = \mu_0\mu_r\vec{H} = \mu\vec{H}$

介质中安培环路定理：$\displaystyle\oint_L \vec{H}\cdot\mathrm{d}\vec{l} = \sum I_0$（传导电流）

\end{document}
